% SEGMENT OLP-0379-S01
% Part: lambda-calculus
% Chapter: lambda-definability
% Section: fixpoints

\documentclass[../../../include/open-logic-section]{subfiles}

\begin{document}

\olfileid{lam}{ldf}{fp}

\olsection{நிலைப்புள்ளிகள்}

பின்வரும் பண்பைக் கொண்ட $\fn{Fac}$ என்ற சொல்லாகக் காரணியச் சார்பை
மீள்வரையறையால் வரையறுக்க விரும்புகிறோம் என்று கொள்க:
\[
\fn{Fac} \ident \lambd[n][\fn{IsZero}\, n\, \num 1 (\fn{Mult}\, n (\fn{Fac}(\fn{Pred} \, n)))]
\]
அதாவது, $n = 0$ எனில் $n$ இன் காரணியம் $1$; இல்லையெனில் அது $n$ மடங்கு
$n-1$ இன் காரணியம். ஆனால் இவ்விதத்தில் $\fn{Fac}$ என்ற சொல்லை வரையறுக்க
முடியாது; ஏனெனில் வலப்பக்கத்திலேயே $\fn{Fac}$ வருகிறது. இத்தகைய
தன்குறிப்பைக் கொண்ட மீள்வரையறைகள் லாம்டா கலனத்தின் பகுதி அல்ல.
% SOURCE CORRECTION TA-LAM-041: use both operands and the Church-zero term in the nonrecursive multiplication example.
எடுத்துக்காட்டாக, ஒரு சொல்லைப் பின்வருமாறு வரையறுப்பது
\[
\fn{Mult} \ident \lambd[ab][a (\fn{Add}\, b) \num{0}]
\]
வலப்பக்கத்தில் $\fn{Add}$ போன்ற முன்பே வரையறுக்கப்பட்ட சொற்களை மட்டுமே
கொண்டுள்ளது. $\fn{Add}$ ஐ அதன் வரையறுக்கும் சொல்லால் பதிலிட்டு எப்போதும்
நீக்கலாம். இதனால் $\fn{Mult}$ என்பது ஒரு தூய லாம்டா சொல்லாகக் கிடைக்கும்.
$\fn{Add}$ தானே வரையறுக்கப்பட்ட சொற்களைக் கொண்டிருந்தால்
($\fn{Add}'$ அவ்வாறு கொண்டிருப்பதுபோல்), இச்செயல்முறையைத் தொடர்ந்து இறுதியில்
ஒரு தூய லாம்டா சொல்லை அடையலாம்.

ஆனால் மேலுள்ள $\fn{Fac}$ வரையறையைப் போன்ற மீள்வரையறைகளுக்கு இது மெய்யல்ல.
வலப்பக்கத்தில் வரும் $\fn{Fac}$ ஐ $\fn{Fac}$ இன் வரையறையாலேயே பதிலிட்டால்,
% SOURCE CORRECTION TA-LAM-042: keep the recursive expansion inside the body of the outer abstraction.
பின்வருவது கிடைக்கிறது:
\begin{align*}
  \fn{Fac} & \ident \lambd[n][\fn{IsZero}\, n\, \num{1} (\fn{Mult}\, n ((\lambd[n][\fn{IsZero} \, n \, \num 1\, (\fn{Mult}\, n\, (\fn{Fac} (\fn{Pred}\, n)))]) (\fn{Pred}\, n)))]
\end{align*}
இப்போதும் வலப்பக்கத்திலிருந்து $\fn{Fac}$ நீங்கவில்லை. இச்செயல்முறையை
மீண்டும் மீண்டும் செய்தால், வரையறை தொடர்ந்து நீள்கிறது; ஒருபோதும் ஒரு தூய
லாம்டா சொல்லாக முடிவதில்லை. ஆகவே இவ்விதத்தில் காரணியத்தை அல்லது பொதுவாக
மீள்வரையறைச் சார்புகளை வரையறுப்பது நடைமுறைக்கு ஒவ்வாதது.

எனினும், மீள்வரையறை நமக்கு ஒன்றைக் கூறுகிறது. காரணியச் சார்பைப்
பிரதிநிதித்துவப்படுத்தும் சொல் $f$ ஒன்று இருக்குமானால், பின்வரும் சொல்
\[
\fn{Fac}' \ident \lambd[g][\lambd[n][\fn{IsZero} \, n \, \num 1\, (\fn{Mult} \, n\, (g (\fn{Pred} n)))]]
\]
$f$ மீது செயல்படுத்தப்படும்போது, அதாவது $\fn{Fac}'\,f$ என்பதும் காரணியச்
சார்பைப் பிரதிநிதித்துவப்படுத்தும். வேறு சொற்களில், ஒரு சார்பை ஏற்று ஒரு
சார்பைத் திருப்பித்தரும் சார்பாக $\fn{Fac}'$ ஐக் கருதினால், $f$ காரணியச்
சார்பாக இருக்கும் நிபந்தனையில் $\fn{Fac'}\, f$ இன் மதிப்பு~$f$ மட்டுமே.
$\fn{Fac}' \, f \equal[\beta] f$ என்ற பண்புள்ள சார்பு~$f$ ஐ
$\fn{Fac}'$ இன் \emph{நிலைப்புள்ளி} என்கிறோம். ஆகவே காரணியச் சார்பு
$\fn{Fac}'$ இன் ஒரு நிலைப்புள்ளி.

கொடுக்கப்பட்ட சொல்லின் நிலைப்புள்ளிகளைக் கணிக்கும் சொற்கள் லாம்டா கலனத்தில்
உள்ளன. அச்சொற்களைக் கொண்டு $\fn{Fac}'$ போன்ற ஒரு சொல்லை காரணியத்தின்
வரையறையாக மாற்றலாம்.

\begin{defn}
  \ollabel{defn:Turing-Y}
  \emph{Y-சேர்ப்பான்} என்பது பின்வரும் சொல்:
  \[
  Y \ident (\lambd[ux][x(uux)])(\lambd[ux][x(uux)]).
  \]
\end{defn}

\begin{thm}
  எந்தச் சொல் $g$ க்கும் $Yg \red g(Yg)$ என்ற பண்பு $Y$ க்கு உள்ளது.
  ஆகவே $Yg$ எப்போதும்~$g$ இன் ஒரு நிலைப்புள்ளி.
\end{thm}

\begin{proof}
  $(\lambd[ux][x(uux)])$ ஐ~$U$ எனச் சுருக்கி எழுதுவோம்; அப்போது
  $Y \ident UU$. எனவே
  \begin{align*}
    Y g &\ident (\lambd[ux][x(uux)])U\,g \\
    &\red (\lambd[x][x(UUx)])g\\
    & \red g(UUg) \ident g(Yg).
  \end{align*}
  $g(Yg)$, $Yg$ இரண்டும் $g(Yg)$ ஆகக் குறைவதால்
  $g(Yg) \equal[\beta] Yg$. ஆகவே $Yg$ என்பது~$g$ இன் ஒரு நிலைப்புள்ளி.
\end{proof}

$Yg$ ஒரு ரெடெக்ஸ் என்பதால், குறைத்தல் முடிவுறாமல் தொடரலாம்:
\begin{align*}
  Y g &\red g (Y g) \\
    &\red g (g (Y g)) \\
    &\red g(g (g (Y g)))\\
    &\ldots
\end{align*}
எனவே $Yg$ என்பதை $g$ தன்னையே முடிவுறாமல் பலமுறை செயல்படுத்துவதாகக்
கருதலாம். அதன் மீது~$g$ ஐ மேலும் ஒருமுறை செயல்படுத்தினால்,
ஒருவகையில் சொல்வதானால், புதிதாக எதையும் செய்வதில்லை. $Yg$ மீது முடிவுறாமல்
பலமுறை செயல்படுத்தப்பட்ட~$g$ மீது மீண்டும்~$g$ ஐச் செயல்படுத்தினாலும்,
அது $Yg$ மீது முடிவுறாமல் பலமுறை செயல்படுத்தப்பட்ட~$g$ ஆகவே உள்ளது.

மேலே~$Yg$ இலிருந்து தொடங்கும் $\beta$-குறைத்தல் படிகளின் தொடர் முடிவுறாதது
என்பதைக் கவனிக்க. ஆகவே $Yg$ ஐ ஏதேனும் ஒரு சொல்லின்மீது செயல்படுத்தி
$(Yg)N$ ஐக் கருதினால், அச்சொல்லும் $(g(Yg))N$, $(g(g(Yg)))N$, \dots\
என முடிவுறாமல் பல வேறு சொற்களாகக் குறையும். இருப்பினும்,
\emph{வேறொரு} குறைத்தல் படித்தொடர் ஓர் இயல்புவடிவில் முடிவடைவது
சாத்தியமே.

காரணியத்தை எடுத்துக்கொள்க. $\fn{Fac}$ ஐ $Y\, \fn{Fac}'$ என, அதாவது
$\fn{Fac'}$ இன் ஒரு நிலைப்புள்ளியாக, வரையறுக்க. அப்போது:
\begin{align*}
  \fn{Fac}\,\num 3
  &\red Y\, \fn{Fac}' \, \num 3 \\
  &\red \fn{Fac}' (Y \,\fn{Fac}') \, \num 3 \\
  &\ident (\lambd[x][\lambd[n][\fn{IsZero} \, n\, \num 1\,
      (\fn{Mult}\, n\, (x (\fn{Pred}\, n)))]]) \, \fn{Fac} \, \num 3\\
  & \red \fn{IsZero} \, \num 3\, \num 1\,
      (\fn{Mult}\, \num 3\, (\fn{Fac} (\fn{Pred}\, \num 3))) \\
  &\red  \fn{Mult}\, \num 3 \, (\fn{Fac} \, \num 2).
  \intertext{இதேபோல்,}
  \fn{Fac}\,\num 2
  &\red  \fn{Mult}\, \num 2 \, (\fn{Fac} \, \num 1) \\
  \fn{Fac}\,\num 1
  &\red  \fn{Mult}\, \num 1 \, (\fn{Fac} \, \num 0)
  \intertext{ஆனால்}
  \fn{Fac}\,\num 0
  &\red \fn{Fac}' (Y \,\fn{Fac}') \, \num 0 \\
  &\ident (\lambd[x][\lambd[n][\fn{IsZero} \, n\, \num 1\,
      (\fn{Mult}\, n\, (x (\fn{Pred}\, n)))]]) \, \fn{Fac} \, \num 0\\
  & \red \fn{IsZero} \, \num 0\, \num 1\,
      (\fn{Mult}\, \num 0\, (\fn{Fac} (\fn{Pred}\, \num 0))).\\
  &\red \num 1.
  \intertext{எனவே எல்லாவற்றையும் சேர்த்து}
  \fn{Fac}\,\num 3
  &\red  \fn{Mult}\, \num 3 \,
  (\fn{Mult} \, \num 2\, (\fn{Mult} \, \num 1\, \num 1)).
\end{align*}

$\fn{Fac'}$ க்குப் பொருந்துவது ஒவ்வொரு மீள்வரையறைக்கும் பொருந்தும்.
பின்வரும் மீள்வரையறைச் சமன்பாடு நம்மிடம் உள்ளதாகக் கொள்க:
\begin{align*}
g\,x_1\dots x_n & \equal[\beta] N
\intertext{இங்கு $N$ ஆனது~$g$, $x_1$, \dots,~$x_n$ ஆகியவற்றைக்
கொண்டிருக்கலாம். அப்போது பின்வருமாறு அமையும்
$G \ident (Y \lambd[g][\lambd[x_1\dots x_n][N]])$ என்ற சொல் எப்போதும்
உள்ளது:}
G\,x_1 \dots x_n & \equal[\beta] \Subst{N}{G}{g}.
\intertext{ஏனெனில் நிலைப்புள்ளித் தேற்றத்தின்படி,}
G \ident (Y \lambd[g][\lambd[x_1\dots x_n][N]]) & \red \lambd[g][\lambd[x_1\dots x_n][N]](Y \lambd[g][\lambd[x_1\dots x_n][N]])\\
& \ident (\lambd[g][\lambd[x_1\dots x_n][N]])\,G
\intertext{எனவே}
G\,x_1 \dots x_n
& \red (\lambd[g][\lambd[x_1\dots x_n][N]])\,G\,x_1\dots x_n\\
& \red (\lambd[x_1\dots x_n][\Subst{N}{G}{g}])\,x_1\dots x_n\\
  & \red \Subst{N}{G}{g}.
\end{align*}

\olref{defn:Turing-Y} இல் உள்ள $Y$ சேர்ப்பானை ஆலன் டியூரிங்
கண்டறிந்தார். அலான்சோ சர்ச் வேறொரு வடிவத்தை முன்மொழிந்தார்; அதை~$Y_C$ என்று
அழைப்போம்:
\[
Y_C \ident \lambd[g][(\lambd[x][g(xx)])(\lambd[x][g(xx)])].
\]
% SOURCE CORRECTION TA-LAM-043: the comparison is about Church's Y_C, not the Turing combinator Y that the preceding theorem proves reduces directly.
சர்ச்சின் சேர்ப்பான் டியூரிங்கின் சேர்ப்பானைவிடச் சற்று வலிமை குறைந்தது:
$Y_Cg \equal[\beta] g(Y_Cg)$; ஆனால் $Y_Cg \bred g(Y_Cg)$ அல்ல.
$V$ என்பது $\lambd[x][g(xx)]$ என்ற சொல்லாக இருக்கட்டும்; அப்போது
$Y_C \ident \lambd[g][VV]$. எனவே
\begin{align*}
  VV & \ident (\lambd[x][g(xx)])V \red g(VV)
  \text{ ஆகவே}\\
  Y_C g & \ident (\lambd[g][VV])g \red VV \red g(VV), \text{ ஆனால் மேலும்}\\
  g(Y_C g) & \ident g((\lambd[g][VV])g) \red g(VV).
\end{align*}
வேறு சொற்களில், $Y_Cg$, $g(Y_Cg)$ இரண்டும் $g(VV)$ என்ற பொதுச் சொல்லுக்குக்
குறைகின்றன; எனவே $Y_Cg \equal[\beta] g(Y_Cg)$. பயன்பாடுகளுக்கு இது
பெரும்பாலும் போதுமானது.

\end{document}
